\documentclass[11pt,a4paper]{article}
\usepackage[utf8]{inputenc}
\usepackage{mathptmx} % Times New Roman Font
\usepackage[margin=1.1in]{geometry}
\usepackage{setspace}
\usepackage{hyperref}
\usepackage{graphicx}

\hypersetup{
    colorlinks=true,
    linkcolor=blue,
    urlcolor=blue
}

\setstretch{1.2}
\setlength{\parindent}{0pt}
\setlength{\parskip}{0.8em}

\begin{document}

\begin{center}
    {\Large \textbf{Author Biographies}} \\[1.2em]
    \textbf{TURRET OS: Multi-Layered Graph and Rule-AST Architecture Rescuing Insider Threat Detection Against Out-of-Distribution Adversaries} \\[1em]
    \textbf{Target Journal:} \textit{Forensic Science International: Digital Investigation} (Elsevier)
\end{center}

\vspace{1.5em}

\textbf{Ali Hamza} is a cybersecurity researcher at the National University of Sciences and Technology (NUST), Islamabad, Pakistan. His research interests focus on system-level security architecture, digital forensic readiness standards (ISO/IEC 27043), W3C-PROV provenance knowledge graphs, out-of-distribution machine learning robustness, and graph neural network fusion. He has led the architectural design and experimental implementation of multi-layered insider threat detection systems.

\vspace{1.5em}

\textbf{Ghulam Mujtaba} is a computer science and cybersecurity researcher at the National University of Computer and Emerging Sciences (FAST-NUCES), Islamabad, Pakistan. His research focuses on graph representation learning, automated grammar-driven rule engines, high-throughput security telemetry ingestion, and enterprise threat intelligence. He has contributed to the graph topological representation and rule-AST evaluation engines for complex system activity analysis.

\end{document}
